\documentclass[9pt]{report}

\usepackage{talks}
\newcommand{\expect}[1]{\mathbb{E}\!\left[ #1 \right]}
\newcommand{\reals}{\mathbb{R}}
\newcommand{\draw}[2]{#1^{(#2)}}
\usepackage{mathpazo}
\usepackage{sourcecodepro}

\begin{document}

\sf \mbox{}
\\[12pt]
\spc{{\LARGE\bfseries \color{MidnightBlue}{Language models for statisticians:}}
\\[8pt]
\spc{\Large\bfseries \color{MidnightBlue}{from $n$-grams to
    transformers to chatbots}}
\\[24pt]
\noindent 
\spc{\large\bfseries \color{MidnightBlue}{Bob Carpenter}}
\\[2pt]
\spc{\small Center for Computational Mathematics}
\\[2pt]
\spc{\small Flatiron Institute}
\vfill 
\noindent 
\spc{\footnotesize July 2023}
\hfill
\includegraphics[width=1.25in]{img/flatiron_logo.png}

\sld{What is a language model?}
\begin{itemize}
\item \myemph{Language} uses a \myemph{finite} number of
  symbols called \myemph{tokens}
  \subit{we assume a finite \myemph{token set} $\mathsf{Tok}$ of size $K$}
\item Tokens may be letters, words,
  sounds, syllables, etc.
  \begin{subitemize}
  \item letters or words traditional (e.g., Shannon 1948)
  \item GPT uses \myemph{sequences of letters} (average 1.5 tokens per
    English word)
  \end{subitemize}
\item Treat language as a \myemph{stochastic process}
  \subit{$Y = Y_1, Y_2, \ldots$ for random variables $Y_n \in T$}
\item Models typically \myemph{autoregressive}, predicting next
  word from previous
  $$
  p_{Y_{n+1} \mid Y_n, Y_{n-1}, \ldots, Y_1}\!\!\left(y_{n + 1} \mid y_n, y_{n-1}, \ldots, y_1\right).
  $$
\end{itemize}

\sld{$N$-gram language models \hfill \normalsize{(Shannon 1948)}}
\begin{itemize}
\item Assume language process is \myemph{order-$N$ Markov}
  \subit{$p(y_{n+1} \mid y_n, y_{n-1}, \ldots, y_1)
    = p(y_{n+1} \mid y_n, y_{n-1}).$}
\item Even \myemph{GPT-4 is Markovian}
  \subit{in the way that a real computer is technically a
    finite-state machine}
\item Limited context allows \myemph{efficient algorithms}
\end{itemize}
    
\sld{Shannon's $N$-gram models}
\begin{itemize}
\item Assume \myemph{training data} $z = z_1, \ldots, z_N$ for $z_n \in T$
\item Use Laplace's \myemph{law of succession}
    $$p(y_{n+1} \mid y_n) = \frac{1 + \textrm{count of } (y_n,
      y_{n+1}) \textrm{ in } z}{N - 1 + K}$$
    where number of tokens is $|\mathsf{Tok}| = K$
  \subit{equivalent to prior $\textrm{Dirichlet}(2)$}
  \end{itemize}
  

\end{document}
